\chapter{Calculations for judgments of learning}
\label{annex:p19-judgments-of-learning}

\begin{annexitem}{annex:p19.jol-gamma}{thm:p19-jol-gamma}{Gamma, comparable pairs and transformations}
\begin{proof}
The two counts are nonnegative. When their sum is positive,
\(-N_+-N_-\leq N_+-N_-\leq N_++N_-\); division proves
the range. The finite set of comparable pairs consists exactly
of the concordant and discordant pairs. Under its uniform law,
the event that the recalled item has the higher rating has
probability \(N_+/(N_++N_-)\). Algebra gives
\[
\frac{1+G}{2}
=\frac{(N_++N_-)+(N_+-N_-)}{2(N_++N_-)}
=\frac{N_+}{N_++N_-}.
\]
This conditions on comparability and makes no independence
assumption about the original item records.

If all outcomes are the same or all judgments are the same,
every pair is excluded and the denominator vanishes.
Conversely, suppose both outcome classes are nonempty and
there are no comparable pairs. Choose \(i_0\) with outcome
zero and \(j_1\) with outcome one. Every outcome-zero item's
judgment must equal \(q_{j_1}\), and every outcome-one item's
judgment must equal \(q_{i_0}\); the chosen cross pair itself
forces those two values to be equal. All judgments are therefore
equal. This proves the exact undefined case, including \(n=1\).

For strictly increasing \(f\), the sign of
\(f(q_i)-f(q_j)\) equals the sign of \(q_i-q_j\), including
zero. Thus every pair retains its concordant, discordant or
excluded status; both counts and gamma are unchanged.
Such a transformation need not preserve the numerical differences
between a judgment and a conditional outcome rate.
The explicit transformation \(f(x)=.99+.01x\) and the
four-item calculation in the next proof demonstrate changes
in calibration, mean bias and squared error while gamma
remains one.
\end{proof}
\end{annexitem}

\begin{annexitem}{annex:p19.jol-calculations}{thm:p19-jol-calculations}{Complete fictional calculations}
\begin{proof}
For the six-item example the recalled items are \(2,4,5,6\)
and the unrecalled items are \(1,3\).
Compare each recalled item with those two unrecalled items:
\[
\begin{array}{c|cc}
\text{Recalled item}&\text{Against item 1}&\text{Against item 3}\\\hline
2&\text{concordant}&\text{discordant}\\
4&\text{concordant}&\text{tied judgment}\\
5&\text{concordant}&\text{concordant}\\
6&\text{concordant}&\text{concordant}
\end{array}
\]
The six pairs within the four recalled items and the one
pair within the two unrecalled items have tied outcomes
and are excluded. Together with the additional tied-judgment
cross pair, these account for eight excluded pairs out of
all fifteen. The remaining counts are \(N_+=6,N_-=1\).
Thus \(G=(6-1)/(6+1)=5/7\), and conditional concordance
is \(6/7\).

The sum of judgments is \(14/5\), so their mean is \(7/15\).
The recall mean is \(4/6=2/3\); subtraction gives \(-1/5\).
The exact-value outcome rates are \(0,1,1/2,1,1\).
Using the empirical weights \(1/6,1/6,2/6,1/6,1/6\),
the squared calibration error is
\[
\frac16\left[
 0+\left(\frac15-1\right)^2
 +2\left(\frac25-\frac12\right)^2
 +\left(\frac45-1\right)^2+0\right]
=\frac{7}{60}.
\]
The single-coordinate squared forecast loss instead uses
each actual binary outcome:
\[
\frac16\left[
0+\frac{16}{25}+\frac4{25}+\frac9{25}+\frac1{25}+0
\right]=\frac15.
\]
The outcome variance remaining in the tied .4 judgment class
has weight \(2/6\) and conditional variance \(1/4\), giving
\(1/12\). Thus the calibration decomposition agrees:
\(7/60+1/12=1/5\).

For the four-item outcome vector \((0,0,1,1)\), there are four
cross-outcome pairs. Both the original reports \((0,0,1,1)\)
and the transformed reports \((.99,.99,1,1)\) make every one
concordant, so gamma is one. The latter reports have mean
.995 against outcome mean .5, giving bias .495.
Their two conditional outcome rates are zero and one,
so squared calibration error is
\[
\frac12(.99-0)^2+\frac12(1-1)^2=.49005.
\]
The original reports had zero error. This is exactly the
strictly increasing transformation specified in the preceding
proof, not a changed outcome record.

Finally pool the two fictional people. Recalled-item judgments
are .9 and .2, and unrecalled-item judgments are .8 and .1.
The .9 judgment forms two concordant pairs. The .2 judgment
forms one concordant and one discordant pair. Thus pooled
gamma is \((3-1)/(3+1)=1/2\), whereas each person's sole
within-person cross-outcome pair is concordant and gives
gamma one. The pooled mean judgment is .5, equal to mean
recall. Since all four judgment values are distinct, their
empirical conditional outcome rates are the individual
zero-one outcomes, and squared calibration error equals
\[
\frac{.8^2+(.9-1)^2+.1^2+(.2-1)^2}{4}
=\frac{1.30}{4}=\frac{13}{40}.
\]
All these calculations describe stipulated records and do
not constitute an estimate from either published experiment.
\end{proof}
\end{annexitem}
